\chapter{Full configuration interaction theory}
\label{chap:fci}

We will start our discussions of many-body methods with what is called
full configuration interaction theory, or just FCI.  The reason for
this is simple. If the effective degrees of freedom which are relevant
for describing a physical system give a sufficiently good reproduction
of the static properties (and later we will add dynamical properties)
of a many-body system, FCI gives us all relevant eigenpairs
(eigenvalues and eigenstates) and allows us to compute other
expectation values than the energy levels.

Although it would have been natural to start with mean-field methods
(FCI is often labelled as a post Hartree-Fock method), we have chosen
a slightly different route. The reason for this is that FCI,
within a restricted effective space, provides us with the exact
eigenvalues and can thus serve as a benchmark for the approximative
methods discussed in the coming chapters, from mean field methods to
the Tamm-Dancoff (TDA) and random-phase (RPA) approximations, variants of many-body
perturbation theory and coupled-cluster theory, Greens' function
approaches and other. We will use FCI to illustrate how these other
methods can be interpreeted in terms of specific truncations/approximations of
FCI.

We start thus with, what this authour likes to claim as the simplest
possible approach to many-body problems, an approach which involves us
being able to set up a many-body Hamiltonian matrix (within an
effective and reduced Hilbert space), and then simply diagonalize. The
overarching principles and the underlying mathematics are
straightforward, although the actual implementations may easily lead
to endless frustrations and sleepless nights. Let us get started.

%----------------------------------------------------------------------
\section{Slater determinants as a basis}
\label{sec:cibasis}

\index{configuration interaction}
The simplest many-body wave function is a product,
\begin{equation}
  \Psi(x_1,x_2,\dots,x_N) \approx
    \phi_1(x_1)\phi_2(x_2)\cdots\phi_N(x_N),
  \label{eq:5-product}
\end{equation}
because we are really only good at thinking about one particle at a time.
Product functions are easy to work with: if the single-particle states are
orthonormal the products are too, and matrix elements factorise into
one-dimensional integrals.  The price is the complete absence of correlations,
which must then be built up by superposing many, many products.  This is the
trade-off that runs through the whole subject -- a compact representation with
difficult integrals, or easy integrals and very many states.

Antisymmetry turns the product into a determinant, and
Chapter~\ref{chap:mbphys} showed that the two elementary properties of
determinants -- a sign change under the interchange of two columns, and
vanishing when two rows coincide -- are exactly the Pauli principle:
\begin{itemize}
  \item no two particles at the same place (two equal columns);
  \item no two particles in the same state (two equal rows).
\end{itemize}
Beyond $N=4$ the coordinate-space determinant becomes unwieldy, which is why
Chapter~\ref{chap:secondquant} replaced it by the occupation representation.
There a Slater determinant is a string of creation operators acting on the
vacuum, each index appearing at most once because
$a^{\dagger}_i a^{\dagger}_i=0$, and antisymmetry is carried by the algebra
rather than by the notation.

We take that formalism as given.  Writing the reference determinant as
\begin{equation}
  \bket{\Phi_0} = \left(\prod_{i\le F}\hat{a}_i^{\dagger}\right)\bket{0},
  \label{eq:5-reference}
\end{equation}
the excited determinants are
\begin{equation}
  \bket{\Phi_i^a} = \hat{a}_a^{\dagger}\hat{a}_i\bket{\Phi_0},
  \qquad
  \bket{\Phi_{ij}^{ab}} =
    \hat{a}_a^{\dagger}\hat{a}_b^{\dagger}\hat{a}_j\hat{a}_i\bket{\Phi_0},
  \label{eq:5-excitations}
\end{equation}
and in general
\begin{equation}
  \bket{\Phi_{ijk\dots}^{abc\dots}} =
    \hat{a}_a^{\dagger}\hat{a}_b^{\dagger}\hat{a}_c^{\dagger}\cdots
    \hat{a}_k\hat{a}_j\hat{a}_i\bket{\Phi_0},
  \label{eq:5-npnh}
\end{equation}
with the hole indices $i,j,k$ below the Fermi level and the particle indices
$a,b,c$ above it, following the convention of
Section~\ref{sec:particlehole}.  We call such a determinant an $np$-$nh$
excitation, or a \emph{configuration}.

%----------------------------------------------------------------------
\section{The configuration-interaction expansion}
\label{sec:ciexpansion}

The exact ground state is an arbitrary superposition of all of these,
\begin{equation}
  \bket{\Psi_0} = C_0\bket{\Phi_0}
    + \sum_{ai}C_i^a\bket{\Phi_i^a}
    + \sum_{abij}C_{ij}^{ab}\bket{\Phi_{ij}^{ab}} + \cdots
    = \left(C_0+\hat{C}\right)\bket{\Phi_0},
  \label{eq:5-ciexpansion}
\end{equation}
where we have introduced the \emph{correlation operator}
\begin{equation}
  \hat{C} = \sum_{ai}C_i^a\hat{a}_a^{\dagger}\hat{a}_i
    + \sum_{abij}C_{ij}^{ab}
      \hat{a}_a^{\dagger}\hat{a}_b^{\dagger}\hat{a}_j\hat{a}_i + \cdots
  \label{eq:5-correlationoperator}
\end{equation}
The normalisation of $\bket{\Psi_0}$ is at our disposal, and $C_0$ is by
hypothesis non-zero, so we may set $C_0=1$ with proportional changes in all
the other coefficients.  With this \emph{intermediate normalisation},
\begin{equation}
  \braket{\Psi_0}{\Phi_0} = \braket{\Phi_0}{\Phi_0} = 1,
  \qquad
  \bket{\Psi_0} = \left(1+\hat{C}\right)\bket{\Phi_0}.
  \label{eq:5-intermediate}
\end{equation}

It is convenient to compress the notation.  Writing $H$ for a set of
$0,1,\dots,n$ hole indices and $P$ for a set of $0,1,\dots,n$ particle
indices,
\begin{equation}
  \bket{\Psi_0} = \sum_{PH}C_H^P\bket{\Phi_H^P}
    = \left(\sum_{PH}C_H^P\hat{A}_H^P\right)\bket{\Phi_0},
  \label{eq:5-compact}
\end{equation}
with $\hat{A}_H^P$ the operator that creates the corresponding excitation.
Unit normalisation reads $\sum_{PH}|C_H^P|^2=1$, and the energy is
\begin{equation}
  E = \bracket{\Psi_0}{\hat{H}}{\Psi_0}
    = \sum_{PP'HH'}C_H^{*P}
      \bracket{\Phi_H^P}{\hat{H}}{\Phi_{H'}^{P'}}C_{H'}^{P'} .
  \label{eq:5-energy}
\end{equation}

If the expansion runs over \emph{all} determinants that can be built from the
chosen single-particle basis, this is \emph{full configuration interaction}.
The only approximation left is the truncation of the single-particle space
itself.

%----------------------------------------------------------------------
\section{The eigenvalue problem}
\label{sec:cieigenvalue}

Equation~(\ref{eq:5-energy}) is solved by diagonalisation, and it is worth
seeing why the diagonalisation is not merely convenient but is the
variational answer.  We minimise the energy subject to normalisation,
\begin{equation}
  \delta\left[\bracket{\Psi_0}{\hat{H}}{\Psi_0}
    - \lambda\braket{\Psi_0}{\Psi_0}\right] = 0,
  \label{eq:5-variation}
\end{equation}
with $\lambda$ a Lagrange multiplier.  Carrying out the variation,
\begin{equation}
  \sum_{P'H'}\Big\{\delta[C_H^{*P}]
    \bracket{\Phi_H^P}{\hat{H}}{\Phi_{H'}^{P'}}C_{H'}^{P'}
  + C_H^{*P}\bracket{\Phi_H^P}{\hat{H}}{\Phi_{H'}^{P'}}\delta[C_{H'}^{P'}]
  - \lambda\big(\delta[C_H^{*P}]C_{H'}^{P'}
    + C_H^{*P}\delta[C_{H'}^{P'}]\big)\Big\} = 0 .
  \label{eq:5-variation2}
\end{equation}
Since $\delta[C_H^{*P}]$ and $\delta[C_{H'}^{P'}]$ are complex conjugates, it
is necessary and sufficient that the quantities multiplying
$\delta[C_H^{*P}]$ vanish, giving
\begin{equation}
  \sum_{P'H'}\bracket{\Phi_H^P}{\hat{H}}{\Phi_{H'}^{P'}}C_{H'}^{P'}
    = \lambda\,C_H^{P}
  \label{eq:5-eigen}
\end{equation}
for every set $PH$.  Multiplying by $C_H^{*P}$ and summing identifies
$\lambda=E$, so the multiplier is the energy and
\begin{equation}
  \boxed{\;
  \sum_{P'H'}\bracket{\Phi_H^P}{\hat{H}-E}{\Phi_{H'}^{P'}}C_{H'}^{P'} = 0
  \;}
  \label{eq:fullci}
\end{equation}
This is the full CI equation.  The same result follows more directly by
writing $(\hat{H}-E)\bket{\Psi_0}=0$ and projecting successively onto every
$\bbra{\Phi_H^P}$.

Equation~(\ref{eq:fullci}) is a Hermitian eigenvalue problem, and the
diagonalisation is a Rayleigh-Ritz variational calculation in the space
spanned by the chosen determinants: the lowest eigenvalue is the best upper
bound to $E_0$ obtainable in that space, and it can only improve as the space
grows.  This is exactly the property that made the Lanczos algorithm of
Section~\ref{sec:lanczos} so well suited to the problem -- the lowest Ritz
value there is the same variational bound, computed in a Krylov space instead
of a space of determinants.

Everything now depends on being able to build and diagonalise the matrix
$\bracket{\Phi_H^P}{\hat{H}}{\Phi_{H'}^{P'}}$.  The matrix elements are
supplied by the Condon-Slater rules of Section~\ref{sec:wickapps}, and the
diagonalisation by the algorithms of Chapter~\ref{chap:linalg}: Householder
and QL for small matrices, Lanczos for large ones.

%----------------------------------------------------------------------
\section{The structure of the Hamiltonian matrix}
\label{sec:cimatrix}

Suppose we have six fermions below the Fermi level, so that at most
$6p$-$6h$ excitations are possible.  Ordering the determinants by excitation
level, the Hamiltonian matrix has the block structure
\begin{equation}
  \begin{array}{c|ccccccc}
     & 0p0h & 1p1h & 2p2h & 3p3h & 4p4h & 5p5h & 6p6h\\
    \hline
    0p0h & \times & \times & \times & 0 & 0 & 0 & 0\\
    1p1h & \times & \times & \times & \times & 0 & 0 & 0\\
    2p2h & \times & \times & \times & \times & \times & 0 & 0\\
    3p3h & 0 & \times & \times & \times & \times & \times & 0\\
    4p4h & 0 & 0 & \times & \times & \times & \times & \times\\
    5p5h & 0 & 0 & 0 & \times & \times & \times & \times\\
    6p6h & 0 & 0 & 0 & 0 & \times & \times & \times
  \end{array}
  \label{eq:5-blockstructure}
\end{equation}
The zeros are not accidental.  A two-body operator connects Slater
determinants differing by at most two single-particle states, which is the
last of the Condon-Slater rules of Section~\ref{sec:wickapps}.  The matrix is
therefore \emph{banded} in the excitation level, with bandwidth two, and
overwhelmingly sparse: for six particles the fraction of non-zero blocks is
already small, and it falls as the space grows.  Chapter~\ref{chap:linalg}
called this the structure that Lanczos exploits; here we see where it comes
from.

\paragraph{In a Hartree-Fock basis.}
If the single-particle states are the solutions of the Hartree-Fock equations
rather than the eigenstates of $\hat{h}_0$, one more set of elements
vanishes.  Brillouin's theorem, met in Section~\ref{sec:wickapps}, states that
\begin{equation}
  \bracket{\Phi_0}{\hat{H}}{\Phi_i^a} = 0,
  \label{eq:5-brillouin}
\end{equation}
so the reference determinant does not couple to the singles at all,
\begin{equation}
  \begin{array}{c|ccccccc}
     & 0p0h & 1p1h & 2p2h & 3p3h & 4p4h & 5p5h & 6p6h\\
    \hline
    0p0h & \tilde{\times} & 0 & \tilde{\times} & 0 & 0 & 0 & 0\\
    1p1h & 0 & \tilde{\times} & \tilde{\times} & \tilde{\times} & 0 & 0 & 0\\
    2p2h & \tilde{\times} & \tilde{\times} & \tilde{\times} & \tilde{\times}
         & \tilde{\times} & 0 & 0\\
    3p3h & 0 & \tilde{\times} & \tilde{\times} & \tilde{\times}
         & \tilde{\times} & \tilde{\times} & 0\\
    4p4h & 0 & 0 & \tilde{\times} & \tilde{\times} & \tilde{\times}
         & \tilde{\times} & \tilde{\times}\\
    5p5h & 0 & 0 & 0 & \tilde{\times} & \tilde{\times} & \tilde{\times}
         & \tilde{\times}\\
    6p6h & 0 & 0 & 0 & 0 & \tilde{\times} & \tilde{\times} & \tilde{\times}
  \end{array}
  \label{eq:5-blockstructureHF}
\end{equation}
where the tildes are a reminder that a unitary transformation of the
single-particle basis has changed every matrix element.  The correlation
energy is then carried entirely by the doubles at leading order, which is the
starting point for both many-body perturbation theory and coupled-cluster
theory.

The programme \texttt{fci.py} builds this structure explicitly.  For the
pairing model with four levels and four particles the full space has
$\binom{8}{4}=70$ determinants, distributed as $1$, $16$, $36$, $16$, $1$ over
the excitation levels $0$ to $4$, and the computed block table is
\begin{equation}
  \begin{array}{c|ccccc}
     & 0p0h & 1p1h & 2p2h & 3p3h & 4p4h\\
    \hline
    0p0h & \times & 0 & \times & 0 & 0\\
    1p1h & 0 & \times & 0 & \times & 0\\
    2p2h & \times & 0 & \times & 0 & \times\\
    3p3h & 0 & \times & 0 & \times & 0\\
    4p4h & 0 & 0 & \times & 0 & \times
  \end{array}
  \label{eq:5-pairingblocks}
\end{equation}
Even the odd blocks are empty here, because the pairing interaction moves
particles two at a time and can never change the number of broken pairs.  The
symmetry has done more work than Brillouin's theorem would.

%----------------------------------------------------------------------
\section{The exponential wall}
\label{sec:exponentialwall}

\index{exponential wall}\index{Hilbert space!dimension}
Full configuration interaction is exact, and that is precisely its problem.
With $N$ particles distributed over $n$ single-particle states the number of
Slater determinants is
\begin{equation}
  \dim\mathcal{H} = \binom{n}{N},
  \label{eq:5-dimension}
\end{equation}
which is the count of bit patterns with $N$ bits set out of $n$, in the
representation of Section~\ref{sec:bitrep}.  This is a combinatorial
expression, but combinatorial growth turns into exponential growth as soon as
the particle number scales with the size of the basis.  At half filling,
$N=n/2$, Stirling's formula gives
\begin{equation}
  \binom{n}{n/2} \sim \frac{2^{n}}{\sqrt{\pi n/2}} \sim \frac{2^{n}}{\sqrt{n}},
  \label{eq:5-stirling}
\end{equation}
so the dimension doubles every time one single-particle state is added.  The
program \texttt{fci.py} tabulates both sides:
\begin{center}
\renewcommand{\arraystretch}{1.15}
\begin{tabular}{rll}
\toprule
$n$ & $\binom{n}{n/2}$ & $2^{n}/\sqrt{n}$\\
\midrule
$10$  & $2.520\times10^{2}$  & $3.238\times10^{2}$\\
$20$  & $1.848\times10^{5}$  & $2.345\times10^{5}$\\
$40$  & $1.378\times10^{11}$ & $1.738\times10^{11}$\\
$80$  & $1.075\times10^{23}$ & $1.352\times10^{23}$\\
$160$ & $9.205\times10^{46}$ & $1.155\times10^{47}$\\
\bottomrule
\end{tabular}
\end{center}
The estimate is accurate to within a factor of order one over five decades of
$n$ and twenty-five orders of magnitude of dimension.

\paragraph{A realistic example.}
Consider oxygen-16 and let the neutrons occupy the four lowest major
oscillator shells, the $0s$, $0p$, $1s0d$ and $1p0f$ shells, which together
hold $2+6+12+20=40$ single-particle states for each species.  The eight
neutrons can be distributed in
\begin{equation}
  \binom{40}{8} = 7.690\times10^{7}
  \label{eq:5-oxygen}
\end{equation}
ways, and since the protons can be distributed independently the full
proton-neutron space has
\begin{equation}
  \binom{40}{8}^2 = 5.914\times10^{15}
  \label{eq:5-oxygen2}
\end{equation}
Slater determinants.  Rotational and parity symmetry reduce this by orders of
magnitude, but adding a single further major shell, or relaxing the
truncation, restores the growth immediately.  Direct diagonalisation by
Householder and QL, Sections~\ref{sec:householder}
and~\ref{sec:qlalgorithm}, is limited to matrices of dimension around $10^5$
because it stores and modifies the whole matrix.  Lanczos,
Section~\ref{sec:lanczos}, needs only the product $\bm{H}\bm{v}$ and reaches
perhaps $10^{10}$ on a large machine.  Beyond that no rearrangement of the
linear algebra helps: the expansion itself must be truncated.

\paragraph{The same wall in the models of Chapter~\ref{chap:systems}.}
Each of the models of the previous chapter meets Eq.~(\ref{eq:5-stirling}) in
its own way.
\begin{itemize}
  \item The \emph{Lipkin model} with $N$ particles in two $N$-fold degenerate
    levels has $\binom{2N}{N}$ determinants in total and $\binom{N}{N/2}\sim
    2^{N}/\sqrt{N}$ in the maximal-quasispin multiplet.  That the SU(2)
    structure of Section~\ref{sec:lipkin} collapses the problem to an
    $(N+1)$-dimensional matrix is a symmetry accident, not a general escape.
  \item The \emph{pairing model} of Section~\ref{sec:pairingmodel} has
    $\binom{2n}{N}$ determinants for $n$ doubly degenerate levels, of which
    only $\binom{n}{N/2}$ have seniority zero.  Again a symmetry -- here the
    conservation of the number of broken pairs -- takes the square root of
    the dimension, and again it is special to the model.
  \item The \emph{Hubbard model} of Section~\ref{sec:hubbard} has a local
    Hilbert space of dimension four, so $L$ sites give $4^{L}$ states in
    total.  At fixed particle number this becomes
    $\binom{L}{N_{\uparrow}}\binom{L}{N_{\downarrow}}$,
    Eq.~(\ref{eq:4-hubbarddim}), which is again exponential at any extensive
    filling.
  \item A \emph{spin chain}, the limit of the Hubbard model at large $U$, has
    no combinatorial disguise at all:
    $\mathcal{H}=\bigotimes_{i=1}^{L}\mathbb{C}^2$ has dimension $2^{L}$ by
    construction.
\end{itemize}
The pattern is worth stating plainly.  Combinatorial counting and
tensor-product structure are two faces of the same growth; the tensor product
is fundamental, the binomial coefficient is what survives after a conservation
law has been imposed.  A conserved quantity can remove a factor, but never the
exponential.

\begin{notebox}
\textbf{Why the wall is not the end of the story.}  The dimension of
$\mathcal{H}$ counts the coefficients an \emph{arbitrary} state would need.
Physical ground states are not arbitrary.  For local Hamiltonians with a gap,
the entanglement entropy across a cut grows with the area of the cut rather
than with the volume it encloses, and states obeying such an area law are
described to high accuracy by a number of parameters polynomial in $L$.  The
Schmidt decomposition of Section~\ref{sec:schmidt} is the one-cut version of
this statement, and the singular values decaying rapidly is exactly what
makes the truncation work.  Applying it at every bond of a chain gives a
matrix product state,
\begin{equation}
  \bket{\psi} = \sum_{s_1\dots s_L}
    A^{s_1}_1A^{s_2}_2\cdots A^{s_L}_L\,\bket{s_1s_2\dots s_L},
  \label{eq:5-mps}
\end{equation}
whose cost is polynomial in the bond dimension $D$; this is the density-matrix
renormalisation group.  The same idea appears elsewhere in this book in
different clothing: a neural-network quantum state replaces the exponentially
many amplitudes by a polynomially parameterised function of the occupation
pattern, and quantum computing replaces them by $n$ qubits whose Hilbert space
is exponential by construction.  Full CI is the reference against which all
of these compressions are measured, which is why we treat it first.
\end{notebox}

%----------------------------------------------------------------------
\section{Full CI for the pairing model}
\label{sec:fcipairing}

The pairing Hamiltonian of Section~\ref{sec:pairingmodel},
\begin{equation}
  \hat{H} = \xi\sum_{p\sigma}(p-1)\hat{a}^{\dagger}_{p\sigma}
    \hat{a}_{p\sigma}
  - \frac{1}{2}g\sum_{pq}\hat{a}^{\dagger}_{p+}\hat{a}^{\dagger}_{p-}
    \hat{a}_{q-}\hat{a}_{q+},
  \label{eq:5-pairingH}
\end{equation}
is the ideal first test.  It is small enough to diagonalise in the full
space, and Chapter~\ref{chap:systems} already gave the exact answer in the
seniority-zero subspace, so the two calculations can be checked against each
other.

Take four doubly degenerate levels and four particles, $\xi=1$.  The full
space has $\binom{8}{4}=70$ determinants, distributed over the excitation
levels as $1$, $16$, $36$, $16$, $1$.  The seniority-zero subspace used in
Chapter~\ref{chap:systems} has only $\binom{4}{2}=6$ states.  Diagonalising
the $70\times70$ matrix at $g=1$ gives $E_0=0.63554847$, identical to the
value in Table~\ref{tab:pairing} to all digits shown.  The determinants with
broken pairs are present in the basis but carry no amplitude in the ground
state, because the interaction cannot create a broken pair from a paired
configuration.  The symmetry that we imposed by hand in Chapter~\ref{chap:systems}
is recovered automatically by the diagonalisation.

This is the general situation.  A symmetry of $\hat{H}$ block-diagonalises the
FCI matrix, and one may either exploit it, gaining a smaller matrix and a
labelled spectrum, or ignore it, gaining a simpler code at the cost of
diagonalising a larger matrix.  Production codes exploit it; the program
accompanying this chapter does not, precisely so that the two answers can be
compared.

%----------------------------------------------------------------------
\section{Full CI for the Hubbard model}
\label{sec:fcihubbard}

The Hubbard model of Section~\ref{sec:hubbard} is the harder test, and the
more instructive one.  We take a ring of $L=6$ sites at half filling,
$N_{\uparrow}=N_{\downarrow}=3$, and work in the momentum basis, where the
hopping term is diagonal,
\begin{equation}
  \hat{H} = \sum_{k\sigma}\varepsilon_k
    \hat{c}^{\dagger}_{k\sigma}\hat{c}_{k\sigma}
  + \frac{U}{L}\sum_{kk'q}
    \hat{c}^{\dagger}_{k+q\uparrow}\hat{c}^{\dagger}_{k'-q\downarrow}
    \hat{c}_{k'\downarrow}\hat{c}_{k\uparrow},
  \qquad
  \varepsilon_k = -2t\cos k,
  \label{eq:5-hubbardmom}
\end{equation}
with $k=2\pi m/L$.  In this basis the non-interacting ground state is a single
determinant -- the Fermi sea -- so the machinery of
Section~\ref{sec:ciexpansion} applies directly with $\bket{\Phi_0}$ the filled
Fermi sea.  The dimension is $\binom{6}{3}^2=400$, with the excitation levels
populated as $1$, $18$, $99$, $164$, $99$, $18$, $1$.

The interesting quantity is how much of the correlation energy survives
truncation as $U/t$ grows:
\begin{center}
\renewcommand{\arraystretch}{1.2}
\begin{tabular}{rlllll}
\toprule
$U/t$ & $E_{\mathrm{ref}}$ & CIS & CID & CISD & FCI\\
\midrule
$1$ & $-6.500000$ & $-6.50000000$ & $-6.59980159$ & $-6.59989778$ & $-6.60115829$\\
$2$ & $-5.000000$ & $-5.00000000$ & $-5.38877703$ & $-5.39026458$ & $-5.40945685$\\
$4$ & $-2.000000$ & $-2.00000000$ & $-3.41192866$ & $-3.43018837$ & $-3.66870618$\\
$8$ & $\phantom{-}4.000000$ & $\phantom{-}2.48338852$ & $-0.36779217$ & $-1.05692183$ & $-2.04813089$\\
\bottomrule
\end{tabular}
\end{center}
The doubles recover $98.66\%$ of the correlation energy at $U/t=1$ but only
$72.22\%$ at $U/t=8$.  Two features deserve comment.

First, the reference does not couple to the singles at all: the interaction
conserves total momentum, and a single excitation changes it, so
$\bracket{\Phi_0}{\hat{H}}{\Phi_i^a}=0$ exactly, for every $U$.  This is
Brillouin's theorem, Eq.~(\ref{eq:5-brillouin}), enforced here by a symmetry
rather than by a variational condition -- the plane-wave determinant is the
Hartree-Fock solution of a translationally invariant problem.  The CIS matrix
is therefore block diagonal, and its lowest eigenvalue is the smaller of the
reference energy and the lowest eigenvalue of the singles block alone.  For
$U/t\le 4$ the reference wins and CIS reproduces it exactly.  At $U/t=8$ it
does not: the singles block, strongly mixed by the large interaction, has an
eigenvalue at $2.483$, below the reference at $4.000$.  The number in the
table is that eigenvalue, and the state it belongs to is orthogonal to
$\bket{\Phi_0}$.  What CIS reports as a ground state is by then not a
correlated version of the reference at all.

Second, the deterioration with $U/t$ is the signature of \emph{strong
correlation}.  At small $U$ the ground state is a single determinant with
small corrections and any method built on that picture works.  At large $U$
the ground state is a superposition of very many determinants with comparable
weight, the reference is a poor starting point, and a truncation at any fixed
excitation level fails.  This is where the tensor-network and quantum-computing
approaches of the later chapters become attractive, and it is why the Hubbard
model at intermediate $U/t$ remains a benchmark problem.

%----------------------------------------------------------------------
\section{Truncating the expansion}
\label{sec:citruncations}

\index{configuration interaction!truncated}
Since the full expansion is out of reach for all but the smallest systems, we
truncate it.  The natural truncation is by excitation level: keep all
determinants up to $n$ particle-hole pairs and discard the rest.  This gives
the familiar hierarchy
\begin{align}
  \bket{\Psi_{\mathrm{CIS}}} &=
    \left(1+\hat{C}_1\right)\bket{\Phi_0},
  \label{eq:5-cis}\\
  \bket{\Psi_{\mathrm{CID}}} &=
    \left(1+\hat{C}_2\right)\bket{\Phi_0},
  \label{eq:5-cid}\\
  \bket{\Psi_{\mathrm{CISD}}} &=
    \left(1+\hat{C}_1+\hat{C}_2\right)\bket{\Phi_0},
  \label{eq:5-cisd}
\end{align}
and so on to CISDT, CISDTQ, until at $\hat{C}_N$ we are back at full CI.
Each of these is again a variational calculation in a subspace, so each
lowest eigenvalue is an upper bound to $E_0$ and the bounds decrease
monotonically as the truncation is relaxed.  The cost, for a basis of $n$
states and $N$ particles, is roughly $\binom{N}{k}\binom{n-N}{k}$ determinants
at excitation level $k$, that is polynomial of degree $2k$ -- which is why
CISD, at $\mathcal{O}(n^6)$, is affordable and CISDTQ, at
$\mathcal{O}(n^{10})$, is not.

For the pairing model with four levels and four particles the truncated
spaces have dimensions $1$, $17$, $53$, $69$, $70$, and the energies are
\begin{center}
\renewcommand{\arraystretch}{1.2}
\begin{tabular}{rlllll}
\toprule
$g$ & $E_{\mathrm{ref}}$ & CIS & CID & CISD & FCI\\
\midrule
$0.25$ & $1.750000$ & $1.75000000$ & $1.73050700$ & $1.73050700$ & $1.73045985$\\
$0.50$ & $1.500000$ & $1.50000000$ & $1.41759645$ & $1.41759645$ & $1.41677428$\\
$1.00$ & $1.000000$ & $1.00000000$ & $0.64890655$ & $0.64890655$ & $0.63554847$\\
$2.00$ & $0.000000$ & $0.00000000$ & $-1.35208670$ & $-1.35208670$ & $-1.48965216$\\
\bottomrule
\end{tabular}
\end{center}
The singles change nothing at all, and CISD coincides with CID: the pairing
interaction moves pairs and cannot break them, so it has no matrix element
between the reference and a $1p$-$1h$ state, nor between a $1p$-$1h$ state and
the reference through any intermediate.  The doubles carry essentially the
whole correlation energy -- $99.8\%$ at $g=0.25$, still $96.3\%$ at $g=1$,
but only $90.8\%$ at $g=2$.  The pattern is the same as in the Hubbard model:
truncation at fixed excitation level works while the reference is good and
degrades as the coupling grows.

%----------------------------------------------------------------------
\section{Size consistency}
\label{sec:sizeconsistency}

\index{size consistency}
There is a defect of truncated CI that is more serious than the loss of
accuracy, because it does not go away as the basis is improved.  Consider two
identical systems $A$ and $B$, infinitely far apart and non-interacting.  The
Hamiltonian is $\hat{H}_A+\hat{H}_B$ and the exact ground state is the product
$\bket{\Psi_A}\bket{\Psi_B}$, so the exact energy must satisfy
\begin{equation}
  E(AB) = E(A) + E(B) = 2E(A).
  \label{eq:5-sizeconsistent}
\end{equation}
A method that reproduces this is called \emph{size consistent}.

Full CI is size consistent, since the product state lies in the full space.
Truncated CI is not.  The reason is immediate: if $A$ and $B$ are each doubly
excited, the combined state is a \emph{quadruple} excitation of the combined
reference, and CID has no room for it.  The product of two doubles is a
quadruple, and truncating at a fixed excitation level is not a
multiplicatively closed operation.

The program \texttt{fci.py} demonstrates this on two non-interacting copies of
a two-level, two-particle pairing system:
\begin{center}
\renewcommand{\arraystretch}{1.2}
\begin{tabular}{rllll}
\toprule
$g$ & $2E(A)$ & $E(AB)$, FCI & $E(AB)$, CID & CID error\\
\midrule
$0.25$ & $-0.26556444$ & $-0.26556444$ & $-0.26550480$ & $5.96\times10^{-5}$\\
$0.50$ & $-0.56155281$ & $-0.56155281$ & $-0.56066017$ & $8.93\times10^{-4}$\\
$1.00$ & $-1.23606798$ & $-1.23606798$ & $-1.22474487$ & $1.13\times10^{-2}$\\
\bottomrule
\end{tabular}
\end{center}
Full CI passes to machine precision; CID fails, and the failure grows with the
coupling.  With $M$ subsystems instead of two the error grows with $M$, so the
energy per particle computed with truncated CI drifts as the system is made
larger.  For an extended system -- nuclear matter, the electron gas, a large
molecule -- this is fatal.

The repair is to make the ansatz multiplicative rather than additive.
Replacing $1+\hat{C}$ by $\exp(\hat{T})$ with $\hat{T}=\hat{T}_1+\hat{T}_2$
gives
\begin{equation}
  \bket{\Psi_{\mathrm{CCSD}}} = e^{\hat{T}_1+\hat{T}_2}\bket{\Phi_0}
  = \left(1+\hat{T}_1+\hat{T}_2
    + \tfrac{1}{2}\hat{T}_2^2 + \cdots\right)\bket{\Phi_0},
  \label{eq:5-ccsd}
\end{equation}
in which the quadruple excitation $\tfrac12\hat{T}_2^2$ appears automatically,
with its coefficient fixed by the doubles rather than determined
independently.  For two non-interacting subsystems
$e^{\hat{T}}=e^{\hat{T}_A}e^{\hat{T}_B}$ and size consistency is exact.  This
is coupled-cluster theory, and Eq.~(\ref{eq:5-ccsd}) is the reason it, rather
than truncated CI, is the workhorse of modern many-body physics.  We shall
develop it properly in a later chapter; the point here is that its central
idea is a response to a defect that full CI lets us see clearly.

%----------------------------------------------------------------------
\section{Truncations and the other many-body methods}
\label{sec:citruncationmap}

We can now do what the introduction to this chapter promised, and read the
standard many-body methods as truncations of Eq.~(\ref{eq:fullci}).  Write the
projected equations explicitly.  Multiplying
$(\hat{H}-E)\bket{\Psi_0}=0$ from the left by $\bbra{\Phi_0}$ and by
$\bbra{\Phi_H^P}$ gives
\begin{align}
  E &= \bracket{\Phi_0}{\hat{H}}{\Phi_0}
    + \sum_{P'H'\ne 0}\bracket{\Phi_0}{\hat{H}}{\Phi_{H'}^{P'}}C_{H'}^{P'},
  \label{eq:5-energyeq}\\
  \left(E - \bracket{\Phi_H^P}{\hat{H}}{\Phi_H^P}\right)C_H^P
   &= \sum_{P'H'\ne PH}
      \bracket{\Phi_H^P}{\hat{H}}{\Phi_{H'}^{P'}}C_{H'}^{P'} .
  \label{eq:5-amplitudeeq}
\end{align}
The first equation says that the correlation energy comes entirely from the
coupling of the reference to the doubles, since only those have a non-zero
matrix element with $\bket{\Phi_0}$.  The second is a set of coupled
equations for the coefficients, which can be solved by iteration: guess the
$C$'s, compute the right-hand side, update, repeat.  Every method in the
remaining chapters is a prescription for which terms of
Eq.~(\ref{eq:5-amplitudeeq}) to keep and how to solve it.

\begin{itemize}
  \item \textbf{Hartree-Fock theory} keeps only $C_0$ and instead varies the
    single-particle basis so that $\bracket{\Phi_0}{\hat{H}}{\Phi_i^a}=0$ --
    the Brillouin condition of Eq.~(\ref{eq:5-brillouin}).  It is the
    optimisation of the reference rather than of the expansion, and it fixes
    the basis in which every subsequent truncation is made.  In the language
    of Eq.~(\ref{eq:5-blockstructureHF}), Hartree-Fock chooses the unitary
    transformation that empties the $0p0h$-$1p1h$ block.
  \item \textbf{The Tamm-Dancoff approximation} restricts the space to the
    $1p$-$1h$ block alone and diagonalises it.  It is CIS applied to excited
    states, and it gives the excitation spectrum of the reference at the
    cheapest non-trivial level.
  \item \textbf{The random-phase approximation} enlarges TDA by allowing the
    ground state to contain $2p$-$2h$ correlations, so that the excitation
    operator carries both $1p$-$1h$ and $1h$-$1p$ pieces.  It is a partial
    summation of the doubles block, performed to all orders in the coupling
    but only for a selected class of terms.
  \item \textbf{Many-body perturbation theory} solves
    Eq.~(\ref{eq:5-amplitudeeq}) by expanding in powers of the interaction
    rather than by iterating to convergence.  Setting $E$ on the left-hand
    side to the reference energy and $C_{H'}^{P'}$ on the right to zero gives
    the first-order coefficients; substituting these back gives the
    second-order energy, and so on.  The order in perturbation theory is the
    number of times the right-hand side has been substituted into itself.
  \item \textbf{Coupled-cluster theory} keeps the same truncation in
    excitation level as CISD but writes the wave function multiplicatively,
    Eq.~(\ref{eq:5-ccsd}).  The equations look like
    Eq.~(\ref{eq:5-amplitudeeq}) with additional nonlinear terms, and it is
    those nonlinear terms that restore size consistency.
\end{itemize}

Seen this way the methods are not a list of unrelated techniques but a set of
answers to one question: given that we cannot keep all of
Eq.~(\ref{eq:5-ciexpansion}), which parts do we keep, and do we keep them
additively or multiplicatively?  Full CI is the fixed point that lets us
measure each answer.  Whenever a model is small enough for
Eq.~(\ref{eq:fullci}) to be solved exactly -- the pairing model, the Lipkin
model, a short Hubbard ring -- we should solve it, and use the result to
calibrate the approximation before applying it to a system where no
calibration is possible.

%----------------------------------------------------------------------
\section{Practical considerations}
\label{sec:cipractical}

Three remarks on implementation, which anticipate the code discussed in the
companion notebook.

\paragraph{Never store the matrix.}
The FCI matrix is sparse, with the band structure of
Eq.~(\ref{eq:5-blockstructure}), and for interesting dimensions it cannot be
stored.  The Lanczos algorithm of Section~\ref{sec:lanczos} needs only the
product $\bm{H}\bm{v}$, and that product is computed by looping over the
determinants and over the non-zero matrix elements of $\hat{H}$, using the bit
representation of Section~\ref{sec:bitrep} to find which determinant each term
produces and with which sign.  This is the single most important structural
choice in an FCI code.

\paragraph{Reorthogonalise.}
The Lanczos vectors lose orthogonality as soon as an eigenvalue has converged,
and spurious copies of converged eigenvalues appear.  Section~\ref{sec:lanczos}
discussed the remedies; for FCI the usual choice is full or selective
reorthogonalisation, at the price of storing the Lanczos vectors.

\paragraph{Use the symmetries.}
Angular momentum, parity, particle number, total momentum and spin projection
all commute with $\hat{H}$ and block-diagonalise the matrix.  Working in a
single symmetry block reduces the dimension by orders of magnitude and, just
as importantly, labels the states.  The pairing calculation of
Section~\ref{sec:fcipairing} illustrated both options -- $70$ determinants
without the symmetry, $6$ with it, the same ground-state energy.

%----------------------------------------------------------------------
\section{Summary}
\label{sec:ch5summary}

Full configuration interaction is the statement that a many-body state is a
superposition of all Slater determinants that the single-particle basis
allows, together with the observation that finding the best superposition is a
Hermitian eigenvalue problem, Eq.~(\ref{eq:fullci}).  Everything else in this
chapter follows from that.  The Condon-Slater rules make the matrix banded in
the excitation level; a Hartree-Fock basis empties one more block; the
combinatorics of Eq.~(\ref{eq:5-dimension}) makes the matrix too large to
build beyond a few tens of single-particle states; and the truncations that
this forces upon us are precisely the approximations that the following
chapters elevate into methods.

The program \texttt{fci.py} in \texttt{BookMaterial/Programs/} carries out
every calculation quoted above: the pairing model in the full
$70$-dimensional space and in the seniority-zero subspace, the block structure
of the Hamiltonian, the CIS/CID/CISD/FCI comparison for pairing and for the
Hubbard ring, the size-consistency test, and the dimension tables of
Section~\ref{sec:exponentialwall}.  The accompanying notebook
\texttt{fci.ipynb} runs the same code cell by cell.

%----------------------------------------------------------------------
\section{Exercises}
\label{sec:ch5exercises}
%=============================================================================

\subsection*{Exercise 1: counting determinants}
\label{ex:5-counting}
\addcontentsline{toc}{subsection}{Exercise 1: counting determinants}

\begin{enumerate}
  \item[a)] Show that the number of Slater determinants with $N$ particles in
  $n$ single-particle states that are $k$ particle-hole excitations of a given
  reference is $\binom{N}{k}\binom{n-N}{k}$, and verify that
  $\sum_k\binom{N}{k}\binom{n-N}{k}=\binom{n}{N}$.
  \item[b)] Evaluate this for $n=8$, $N=4$ and compare with the distribution
  $1,16,36,16,1$ quoted in Section~\ref{sec:cimatrix}.
  \item[c)] Use Stirling's formula to derive Eq.~(\ref{eq:5-stirling}), and
  determine the largest $n$ for which $\binom{n}{n/2}$ stays below $10^{10}$.
\end{enumerate}

\paragraph{Answer to a).}
Choosing which $k$ of the $N$ occupied states to empty can be done in
$\binom{N}{k}$ ways, and which $k$ of the $n-N$ empty states to fill in
$\binom{n-N}{k}$ ways; the two choices are independent.  Every determinant
with $N$ particles arises in this way exactly once, for exactly one value of
$k$, namely the number of occupied states it does not share with the
reference.  Summing over $k$ must therefore give the total, which is
Vandermonde's identity
\begin{equation}
  \sum_k\binom{N}{k}\binom{n-N}{k}
  = \sum_k\binom{N}{N-k}\binom{n-N}{k} = \binom{n}{N}.
  \label{eq:5-exa1}
\end{equation}

\paragraph{Answer to b).}
With $n=8$ and $N=4$ the products are $\binom{4}{k}^2 = 1,16,36,16,1$ for
$k=0,\dots,4$, summing to $70=\binom{8}{4}$.  This is exactly the distribution
that \texttt{fci.py} finds by classifying the $70$ determinants of the pairing
model by excitation level.

\paragraph{Answer to c).}
Stirling's formula $m!\sim\sqrt{2\pi m}\,(m/e)^m$ gives
\begin{equation}
  \binom{n}{n/2}
  = \frac{n!}{[(n/2)!]^2}
  \sim \frac{\sqrt{2\pi n}\,(n/e)^n}
            {2\pi (n/2)\,\left[(n/2e)^{n/2}\right]^2}
  = \frac{2^{n}}{\sqrt{\pi n/2}} ,
  \label{eq:5-exa2}
\end{equation}
which is Eq.~(\ref{eq:5-stirling}) up to the constant $\sqrt{\pi/2}\approx
1.25$.  Solving $\binom{n}{n/2}=10^{10}$ numerically: $\binom{36}{18}=
9.08\times10^{9}$ and $\binom{38}{19}=3.49\times10^{10}$, so $n=36$ is the
last even value below the limit.  Every further pair of single-particle
states multiplies the dimension by about four.

\subsection*{Exercise 2: which blocks vanish}
\label{ex:5-blocks}
\addcontentsline{toc}{subsection}{Exercise 2: which blocks vanish}

Let $\hat{H}=\hat{H}_0+\hat{H}_I$ with $\hat{H}_I$ a two-body operator.

\begin{enumerate}
  \item[a)] Show that $\bracket{\Phi_H^P}{\hat{H}}{\Phi_{H'}^{P'}}=0$ whenever
  the two determinants differ by more than two single-particle states, and
  deduce the band structure of Eq.~(\ref{eq:5-blockstructure}).
  \item[b)] Suppose $\hat{H}$ also contains a genuine three-body force.  How
  does the band structure change, and what does that do to the cost of a
  matrix-vector product?
  \item[c)] For the pairing Hamiltonian, Eq.~(\ref{eq:5-pairingH}), show that
  every block connecting an even to an odd excitation level vanishes, as in
  Eq.~(\ref{eq:5-pairingblocks}).
\end{enumerate}

\paragraph{Answer to a).}
Write both determinants in the occupation representation of
Section~\ref{sec:bitrep}.  A two-body operator is a product of two creation
and two annihilation operators, so it can change at most two occupation
numbers from one to zero and two from zero to one.  If the two determinants
differ in more than four occupation numbers -- that is, by more than two
single-particle states -- then in the Wick expansion of
Section~\ref{sec:wickapps} at least one operator is left uncontracted in every
term, and the matrix element is zero.  Since an $np$-$nh$ and an $mp$-$mh$
determinant differ by at least $|n-m|$ single-particle states, the element
vanishes for $|n-m|>2$.  The matrix is therefore banded with bandwidth two in
the excitation level, which is Eq.~(\ref{eq:5-blockstructure}).

\paragraph{Answer to b).}
A three-body operator changes up to three occupations of each kind, so the
bandwidth becomes three and the $0p0h$ row now reaches the $3p3h$ block.  The
number of non-zero elements per row grows from $\mathcal{O}(N^2(n-N)^2)$ to
$\mathcal{O}(N^3(n-N)^3)$, so a matrix-vector product costs a further factor
$N(n-N)$ -- typically two to three orders of magnitude.  This is why
three-body forces, though physically important in nuclear physics, are usually
normal-ordered with respect to the reference determinant and then truncated at
the two-body level, a procedure whose justification rests on
Section~\ref{sec:normalordering}.

\paragraph{Answer to c).}
The interaction
$-\tfrac12 g\sum_{pq}\hat{a}^{\dagger}_{p+}\hat{a}^{\dagger}_{p-}
\hat{a}_{q-}\hat{a}_{q+}$ removes a complete pair from level $q$ and creates a
complete pair in level $p$.  Define the seniority $s$ as the number of
unpaired particles.  Both the one-body term and the interaction leave $s$
unchanged, so $[\hat{H},\hat{s}]=0$ and $\hat{H}$ is block diagonal in
seniority.  An excitation of odd order relative to the paired reference must
leave at least two particles unpaired, so it has $s\ge2$, whereas the
reference and all even excitations reached from it by moving whole pairs have
$s=0$.  Every even-odd matrix element therefore connects different seniorities
and vanishes.  Numerically the $70\times70$ matrix splits into a
six-dimensional $s=0$ block, a $s=2$ block and so on, and the ground state
lies in the first.

\subsection*{Exercise 3: full CI for the pairing model}
\label{ex:5-pairingfci}
\addcontentsline{toc}{subsection}{Exercise 3: full CI for the pairing model}

Take four doubly degenerate levels with single-particle energies
$\varepsilon_p=(p-1)\xi$, $p=1,\dots,4$, four particles, $\xi=1$, and the
Hamiltonian of Eq.~(\ref{eq:5-pairingH}).

\begin{enumerate}
  \item[a)] Set up the six seniority-zero configurations and write down the
  $6\times6$ Hamiltonian matrix.
  \item[b)] Diagonalise it for $g=-1,-0.5,0,0.5,1$ and reproduce
  Table~\ref{tab:pairing}.
  \item[c)] Repeat the calculation in the full $70$-dimensional space and
  confirm that the ground-state energy is unchanged.  Inspect the eigenvector
  and verify that every determinant with a broken pair has zero amplitude.
  \item[d)] Plot the correlation energy against $g$ and show that it is
  quadratic for small $|g|$, as second-order perturbation theory requires.
\end{enumerate}

\paragraph{Answer to a).}
Label a seniority-zero configuration by the pair of levels that are occupied.
With four levels and two pairs there are $\binom{4}{2}=6$ of them,
\begin{equation}
  \bket{12},\;\bket{13},\;\bket{14},\;\bket{23},\;\bket{24},\;\bket{34},
  \label{eq:5-exb1}
\end{equation}
where $\bket{pq}$ has pairs in levels $p$ and $q$.  The diagonal element of
$\bket{pq}$ is $2\xi[(p-1)+(q-1)]-g$: the factor two because each level holds
two particles, and $-g$ because the $p=q$ terms of the interaction give
$-\tfrac12 g$ for each of the two occupied pairs.  A single application of the
interaction moves \emph{one} pair, so two configurations are connected by
$-\tfrac12 g$ if and only if they share exactly one occupied level; the three
pairs of configurations that share no level, $\bket{12}$--$\bket{34}$,
$\bket{13}$--$\bket{24}$ and $\bket{14}$--$\bket{23}$, are not connected at
all.  With $\xi=1$ and the ordering of Eq.~(\ref{eq:5-exb1}),
\begin{equation}
  \bm{H} =
  \begin{pmatrix}
    2-g & -g/2 & -g/2 & -g/2 & -g/2 & 0\\
    -g/2 & 4-g & -g/2 & -g/2 & 0 & -g/2\\
    -g/2 & -g/2 & 6-g & 0 & -g/2 & -g/2\\
    -g/2 & -g/2 & 0 & 6-g & -g/2 & -g/2\\
    -g/2 & 0 & -g/2 & -g/2 & 8-g & -g/2\\
    0 & -g/2 & -g/2 & -g/2 & -g/2 & 10-g
  \end{pmatrix}.
  \label{eq:5-exb2}
\end{equation}

\paragraph{Answer to b).}
Diagonalising Eq.~(\ref{eq:5-exb2}) reproduces Table~\ref{tab:pairing}; at
$g=1$ the lowest eigenvalue is $0.63554847$ and the reference energy is
$E_{\mathrm{ref}}=H_{11}=1$.

\paragraph{Answer to c).}
The class \texttt{PairingFCI} in \texttt{fci.py} builds the full space with no
seniority restriction.  Its dimension is $\binom{8}{4}=70$, the excitation
levels are populated as $1,16,36,16,1$, and the lowest eigenvalue at $g=1$ is
$0.63554847$, identical to b) to every digit.  Projecting the ground-state
eigenvector onto the $64$ determinants with at least one broken pair gives
zero to machine precision, as Exercise 2c) predicts.

\paragraph{Answer to d).}
The reference is $\bket{12}$.  Second-order perturbation theory in the
interaction gives
\begin{equation}
  E^{(2)} = \sum_{\Phi\ne\Phi_0}
    \frac{|\bracket{\Phi_0}{\hat{H}_I}{\Phi}|^2}{E_0^{(0)}-E_{\Phi}^{(0)}}
  = \frac{g^2}{4}\left(\frac{1}{-2}+\frac{1}{-4}+\frac{1}{-4}
    +\frac{1}{-6}\right)
  = -\frac{7}{24}g^2 ,
  \label{eq:5-exb3}
\end{equation}
using the four configurations reachable from $\bket{12}$ by moving one pair,
with unperturbed excitation energies $2,4,4,6$.  A log-log plot of
$|E_{\mathrm{corr}}|$ against $g$ has slope $2$ for $|g|\lesssim0.2$ and bends
away above that, where the higher orders take over.  At $g=0.25$,
Eq.~(\ref{eq:5-exb3}) gives $-0.01823$ against the exact $-0.01954$; at
$g=0.5$ it gives $-0.0729$ against $-0.0832$, and at $g=1$ only $-0.2917$
against $-0.3645$.

\subsection*{Exercise 4: truncations and size consistency}
\label{ex:5-sizeconsistency}
\addcontentsline{toc}{subsection}{Exercise 4: truncations and size consistency}

\begin{enumerate}
  \item[a)] Using the code of the previous exercise, restrict the basis to
  excitation levels $\le1$ and $\le2$ and reproduce the CIS and CID columns of
  Section~\ref{sec:citruncations}.  Explain why CIS returns exactly the
  reference energy.
  \item[b)] Build a system of two identical, non-interacting two-level pairing
  systems.  Show that full CI gives exactly twice the energy of one subsystem
  while CID does not, and tabulate the error for $g=0.25,0.5,1$.
  \item[c)] Identify the configuration that CID is missing, and show that
  including quadruple excitations restores the exact result.
  \item[d)] Show that the ansatz $\exp(\hat{T}_2)\bket{\Phi_0}$, with the same
  number of independent parameters as CID, is size consistent for this system.
\end{enumerate}

\paragraph{Answer to a).}
In a truncated space the eigenvalue problem is Eq.~(\ref{eq:fullci})
restricted to the corresponding rows and columns.  The dimensions for the cuts
at $0,1,2,3,4$ are $1,17,53,69,70$.  CIS returns the reference energy because
$\bracket{\Phi_0}{\hat{H}}{\Phi_i^a}=0$: a single excitation breaks a pair and
so changes the seniority, which the pairing Hamiltonian conserves
(Exercise 2c).  The reference is therefore decoupled from the
entire singles block and remains an exact eigenstate of the truncated matrix.
Its energy is also the lowest diagonal element, and the singles block lies far
above it for the couplings considered, so it is the CIS ground state.

\paragraph{Answer to b).}
Let each subsystem have two levels at $0$ and $\xi$, holding one pair.  The
combined system then has four levels with energies $0,\xi,0,\xi$ and two
pairs, one per subsystem, with the interaction acting only within a
subsystem.  The class \texttt{NonInteractingCopies} gives
\begin{equation}
\begin{array}{rllll}
g & 2E(A) & E(AB)_{\mathrm{FCI}} & E(AB)_{\mathrm{CID}}
  & \text{CID error}\\[2pt]
0.25 & -0.26556444 & -0.26556444 & -0.26550480 & 5.96\times10^{-5}\\
0.50 & -0.56155281 & -0.56155281 & -0.56066017 & 8.93\times10^{-4}\\
1.00 & -1.23606798 & -1.23606798 & -1.22474487 & 1.13\times10^{-2}
\end{array}
\label{eq:5-exc1}
\end{equation}
Full CI is exact to machine precision; the CID error grows by roughly an order
of magnitude each time $g$ is doubled.

\paragraph{Answer to c).}
Writing $\hat{A}_2^X$ for the operator that excites the pair of subsystem $X$,
the exact product state is
\begin{equation}
  \left(1+c\hat{A}_2^A\right)\left(1+c\hat{A}_2^B\right)\bket{\Phi_0}
  = \left(1 + c\hat{A}_2^A + c\hat{A}_2^B
    + c^2\hat{A}_2^A\hat{A}_2^B\right)\bket{\Phi_0}.
  \label{eq:5-exc2}
\end{equation}
The last term excites both subsystems at once and is a $4p$-$4h$ excitation of
the combined reference; it is absent from CID.  Adding the quadruples -- which
for this small system means going to the full space -- recovers the exact
energy, and the coefficient of the quadruple comes out equal to $c^2$, the
product of the two double amplitudes.

\paragraph{Answer to d).}
With $\hat{T}_2 = t\left(\hat{A}_2^A+\hat{A}_2^B\right)$ the two terms commute,
since they act on disjoint sets of orbitals, so
\begin{equation}
  e^{\hat{T}_2} = e^{t\hat{A}_2^A}e^{t\hat{A}_2^B}
  = \left(1+t\hat{A}_2^A\right)\left(1+t\hat{A}_2^B\right),
  \label{eq:5-exc3}
\end{equation}
the last step because $(\hat{A}_2^X)^2=0$ -- one cannot excite the same pair
twice.  The exponential ansatz therefore generates exactly the product state
of Eq.~(\ref{eq:5-exc2}) with one parameter per subsystem, the same count as
CID, and the energy is additive by construction.  This is coupled cluster with
doubles in miniature, and it shows that the cure for size inconsistency costs
nothing in parameters -- only a change from a linear to a multiplicative
ansatz.

\subsection*{Exercise 5: the Hubbard ring by full CI}
\label{ex:5-hubbardfci}
\addcontentsline{toc}{subsection}{Exercise 5: the Hubbard ring by full CI}

Take a ring of $L=6$ sites at half filling with $t=1$, in the momentum basis
of Eq.~(\ref{eq:5-hubbardmom}).

\begin{enumerate}
  \item[a)] Verify that the dimension is $\binom{6}{3}^2=400$ and that the
  excitation levels are populated as $1$, $18$, $99$, $164$, $99$, $18$, $1$.
  \item[b)] Compute the exact ground-state energy for $U/t=1,2,4,8$ and
  compare with the CID result.
  \item[c)] Show that $\bracket{\Phi_0}{\hat{H}}{\Phi_i^a}=0$ for every single
  excitation, and explain this using conservation of total momentum.
  \item[d)] At $U/t=8$ the CIS energy lies \emph{below} the reference energy
  even though the reference does not couple to the singles.  Explain how this
  is possible and what it says about the reliability of a variational bound.
  \item[e)] Compare the exact energies at large $U$ with the Heisenberg limit
  of Section~\ref{sec:hubbard}.
\end{enumerate}

\paragraph{Answer to a).}
There are $\binom{6}{3}=20$ ways to place three electrons of each spin, and
the two species are independent, so the dimension is $400$.  Counting for each
determinant the number of occupied momenta outside the Fermi sea, separately
for each spin, and adding the two counts gives $1,18,99,164,99,18,1$, which
sums to $400$.  The distribution is symmetric because a $kp$-$kh$ excitation
of the Fermi sea is a $(6-k)p$-$(6-k)h$ excitation of the completely inverted
determinant.

\paragraph{Answer to b).}
The values are those of Section~\ref{sec:fcihubbard}.  Relative to the
reference, the doubles recover $98.66\%$, $94.95\%$, $84.61\%$ and $72.22\%$
of the correlation energy at $U/t=1,2,4,8$.

\paragraph{Answer to c).}
The interaction is
$\tfrac{U}{L}\sum_{kk'q}\hat{c}^{\dagger}_{k+q\uparrow}
\hat{c}^{\dagger}_{k'-q\downarrow}\hat{c}_{k'\downarrow}\hat{c}_{k\uparrow}$,
in which the momenta of the two created operators sum to those of the two
annihilated ones; total crystal momentum is therefore conserved.  A single
excitation $i\to a$ changes the total momentum by $k_a-k_i\ne0$ and so cannot
be reached from the Fermi sea.  The one-body part is diagonal in $k$ and
cannot connect them either.  Evaluating the couplings numerically gives zero
to machine precision for every $U$.  Note that this is Brillouin's theorem,
Eq.~(\ref{eq:5-brillouin}), obtained here from a symmetry rather than from a
variational condition: for a translationally invariant problem the plane-wave
determinant \emph{is} the Hartree-Fock solution.

\paragraph{Answer to d).}
Because the coupling vanishes, the CIS matrix is block diagonal: a $1\times1$
block holding $E_{\mathrm{ref}}$ and an $18\times18$ singles block.  Its
lowest eigenvalue is the smaller of the two numbers.  At $U/t\le4$ the
reference wins and CIS reproduces it exactly.  At $U/t=8$ the singles block,
strongly mixed by the large interaction, has an eigenvalue at $2.483$, which
lies below $E_{\mathrm{ref}}=4.000$, and the variational principle selects the
corresponding state -- one that is exactly orthogonal to $\bket{\Phi_0}$.  The
number is a legitimate upper bound to $E_0$, but it is not a correlated
correction to the reference; the calculation has silently changed which state
it describes.  The moral is that a lower variational energy does not by itself
mean a better description of the state one set out to compute, and that the
overlap with the reference should always be monitored alongside the energy.

\paragraph{Answer to e).}
For $U/t\gg1$ the half-filled Hubbard model maps onto a Heisenberg chain with
$J=4t^2/U$, Eq.~(\ref{eq:4-tJ}), whose ground-state energy per site tends to
$-J(\ln 2-\tfrac14)$ in the thermodynamic limit.  At $U/t=8$ and $L=6$ this
estimate gives about $-1.3$ against the exact $-2.05$.  The discrepancy is
expected: the mapping neglects terms of order $t^4/U^3$, the ring is short, and
$U/t=8$ is not yet deep in the strong-coupling regime.  Repeating the
comparison at $U/t=20$ and $40$ shows the two converging.

\subsection*{Exercise 6: the exponential wall in practice}
\label{ex:5-wall}
\addcontentsline{toc}{subsection}{Exercise 6: the exponential wall in practice}

\begin{enumerate}
  \item[a)] A shell-model code stores the Hamiltonian matrix in double
  precision.  For a dimension of $10^6$ with on average $10^3$ non-zero
  elements per row, and a $4$-byte column index per element, how much memory
  does the sparse matrix need?  Compare with storing the full matrix.
  \item[b)] How large a single-particle space can be treated at half filling
  if a Lanczos code can handle $10^{10}$ determinants?
  \item[c)] A matrix product state with bond dimension $D$ on $L$ sites has
  about $LD^2d$ parameters, with $d$ the local dimension.  For a spin chain
  with $d=2$, $L=100$ and $D=200$, compare this with $2^{L}$.
  \item[d)] Discuss what must be true of the physical state for the comparison
  in c) to be a fair one.
\end{enumerate}

\paragraph{Answer to a).}
The sparse matrix holds $10^{6}\times10^{3}=10^{9}$ elements at $8+4=12$ bytes
each, or $12$~GB -- large but routine.  The full matrix would need
$8\times10^{12}$ bytes, $8$~TB, and diagonalising it by Householder and QL
would take $\mathcal{O}(10^{18})$ floating-point operations.  This is the
practical content of Section~\ref{sec:cipractical}: the dense algorithms of
Sections~\ref{sec:householder} and~\ref{sec:qlalgorithm} are unusable here,
and even the sparse matrix is best not stored at all but regenerated on the
fly inside the matrix-vector product.

\paragraph{Answer to b).}
From Eq.~(\ref{eq:5-stirling}), $2^{n}/\sqrt{n}=10^{10}$ gives $n\approx38$,
and the exact condition $\binom{n}{n/2}\le10^{10}$ gives $n=36$.  A Lanczos
calculation at the limit of current hardware therefore reaches roughly
$36$--$40$ single-particle states at half filling.  Away from half filling the
binomial coefficient is much smaller and the reach is greater, which is why
calculations near closed shells are so much easier than calculations in the
middle of a shell.

\paragraph{Answer to c).}
The matrix product state has about $100\times200^2\times2=8\times10^{6}$
parameters, against $2^{100}=1.27\times10^{30}$ amplitudes for a general
state -- a compression by twenty-four orders of magnitude.

\paragraph{Answer to d).}
The comparison is fair only if the state of interest is actually representable
at that bond dimension.  The Schmidt spectrum across every bond must decay
fast enough that keeping $D$ singular values costs little, which is the
truncation error of the Eckart-Young theorem in
Section~\ref{sec:schmidt} applied bond by bond.  This holds for the ground
states of gapped one-dimensional local Hamiltonians, which obey an area law.
For critical chains the required $D$ grows polynomially with $L$; for
volume-law states -- generic excited states, or the state produced by a
quantum quench after a time proportional to $L$ -- no polynomial $D$
suffices and the exponential wall stands undiminished.  The same caveat
applies to every other compressed representation, including the
neural-network and quantum-circuit ansaetze of the later chapters: each buys
its efficiency by restricting the class of states it can describe, and full
configuration interaction remains the only method that makes no such
restriction at all.
